\documentclass[12pt,a4paper]{article}
\usepackage[margin=2.5cm]{geometry}
\usepackage{booktabs}
\usepackage{hyperref}
\usepackage{parskip}
\usepackage{titlesec}
\usepackage{xcolor}

\definecolor{draftgray}{gray}{0.45}

\title{\textbf{Methodology Draft Notes}\\
\large Coding Agent Adoption in JavaScript/TypeScript/npm Repositories on GitHub\\
\large Geographic Country Ranking Study}
\author{Manuel Youssef Haik Kevorkian}
\date{October 2026 \quad\textcolor{red}{\textbf{DRAFT -- not for submission}}}

\begin{document}
\maketitle
\tableofcontents
\newpage

% ============================================================
\section{Research Questions}
% ============================================================

This study investigates three research questions:

\begin{itemize}
  \item \textbf{RQ1:} What proportion of active JavaScript/TypeScript/npm repositories
        on GitHub show evidence of AI coding agent use, and is this proportion
        statistically non-trivial?

  \item \textbf{RQ2:} How much does single-signal detection undercount AI adoption
        compared to a multi-signal approach, and is this difference statistically
        significant?

  \item \textbf{RQ3:} Do AI coding agent adoption rates differ significantly across
        countries within the JS/TS/npm ecosystem, and which countries show the highest
        and lowest rates?
\end{itemize}

RQ1 and RQ2 address global adoption and measurement reliability. RQ3 introduces the
geographic dimension by ranking all countries with enough repositories in the corpus.

% ============================================================
\section{Pipeline Overview}
% ============================================================

The study follows four sequential phases. No sampling is applied; the full corpus of
60,129 repositories is processed end-to-end.

\begin{enumerate}
  \item \textbf{Phase 1} -- Corpus construction: 60,129 JS/TS/npm repos collected and
        categorized by ecosystem (npm, pnpm, yarn, bun, deno, vanilla).
  \item \textbf{Phase 2} -- Owner location inference: two-pass pipeline assigns a country
        to each of the 50,002 unique owners using GitHub profile signals.
  \item \textbf{Phase 3} -- AI adoption detection: three-pass layered pipeline classifies
        each of the 60,129 repos as agentic, ai\_assisted, or none.
  \item \textbf{Phase 4} -- Geographic analysis: repos joined with owner countries,
        adoption rates computed per country, formal statistical tests applied.
\end{enumerate}

% ============================================================
\section{Phase 1 -- Corpus Construction}
% ============================================================

\subsection{Collection}

Repositories collected via GitHub GraphQL Search API in two passes (JavaScript and
TypeScript separately -- GitHub does not support OR-combined language filters in one
query). Rolling 5-day creation date windows from 2025-02-24 to October 2026. Windows
exceeding the 1,000-result cap were recursively halved until every window fit.

\medskip
\noindent\textbf{Inclusion filters (agreed across research team):}

\begin{table}[h]
\centering
\begin{tabular}{@{}lll@{}}
\toprule
Filter & Value & Rationale \\
\midrule
Language            & JavaScript or TypeScript & Thesis scope \\
Created on or after & 2025-02-24               & Mature AI agent era \\
Stars               & $\geq 4$                 & Excludes toy/abandoned projects \\
Commits             & $\geq 50$                & Meaningful development history \\
Fork                & No                       & Excludes derivative projects \\
JS + TS bytes       & $\geq 175{,}000$         & JS/TS is primary language \\
\bottomrule
\end{tabular}
\end{table}

\noindent\textbf{Result: 60,129 repositories.}

\subsection{Ecosystem Categorization}

Seven files checked at root per repo via GraphQL batch queries (15 repos per request).
Classification priority: deno.json $\to$ deno; package.json + package-lock.json $\to$ npm;
+ yarn.lock $\to$ yarn; + pnpm-lock.yaml $\to$ pnpm; + bun.lockb $\to$ bun;
package.json only $\to$ npm\_no\_lock; none $\to$ vanilla.

\begin{table}[h]
\centering
\begin{tabular}{@{}lrr@{}}
\toprule
Ecosystem & Count & \% \\
\midrule
npm           & 28,248 & 47.0\% \\
pnpm          & 12,229 & 20.3\% \\
vanilla       & 10,177 & 16.9\% \\
npm\_no\_lock &  7,979 & 13.3\% \\
yarn          &  1,091 &  1.8\% \\
deno          &    250 &  0.4\% \\
bun           &    107 &  0.2\% \\
unknown       &     48 &  0.1\% \\
\textbf{Total} & \textbf{60,129} & \textbf{100\%} \\
\bottomrule
\end{tabular}
\end{table}

\noindent\textbf{Note on sampling:} A previous version of this study drew a stratified
sample of 4,500 repos. Sampling has been removed because the layered detection pipeline
(Phase 3) processes the full corpus efficiently without requiring sampling, and sampling
before geographic labeling would introduce bias into the country-level analysis.

\subsection{Output Files}
\begin{itemize}
  \item \texttt{data/stage1/global\_corpus.csv} -- all 60,129 repos
  \item \texttt{data/stage1/ecosystem\_audit.csv} -- category breakdown
\end{itemize}

% ============================================================
\section{Phase 2 -- Owner Location Inference}
% ============================================================

\subsection{Motivation}

GitHub location is attached to user profiles, not repositories. For each of the
50,002 unique owner logins in the corpus, a country is inferred in two passes.
Approximately 60\% of GitHub users leave the location field blank -- this is why
a two-pass approach is needed rather than a simple field lookup.

\subsection{Pass 1 -- Fast Location Lookup}

Batch GraphQL queries (50 owners per request) fetch the raw \texttt{location} profile
field. Each string is matched against a keyword lookup table of country names, major
cities, and common abbreviations in English and Arabic, using substring matching
(longest key first to avoid partial matches).

\medskip
\noindent\textbf{Result: approximately 13,625 owners located (27.2\%). 36,377 passed to Pass 2.}

\subsection{Pass 2 -- Deep Profile Inference}

One GraphQL query per unresolved owner fetches additional profile fields. Six signals
are tried in order; inference stops at the first match:

\begin{table}[h]
\centering
\begin{tabular}{@{}clll@{}}
\toprule
Priority & Signal & Field & Confidence \\
\midrule
1 & Bio keyword match      & \texttt{bio}                  & Medium \\
2 & Company keyword match  & \texttt{company}              & Medium \\
3 & Company domain TLD     & \texttt{company} (domain)     & Medium \\
4 & Email domain TLD       & \texttt{email}                & Medium \\
5 & Website URL TLD        & \texttt{websiteUrl}           & Low \\
6 & Org location keyword   & \texttt{organizations.location} & Medium \\
7 & Network majority       & followers + following ($\geq$3 located, $\geq$40\% agree) & Low \\
\bottomrule
\end{tabular}
\end{table}

Country-code TLDs: .de (Germany), .eg (Egypt), .br (Brazil), .in (India), etc.
Owners unresolved after all signals are labelled \texttt{not\_specified} and excluded
from country-level analysis.

\subsection{Output Files}
\begin{itemize}
  \item \texttt{data/stage2/owner\_locations\_fast.csv} -- Pass 1 results
  \item \texttt{data/stage2/owner\_locations\_deep.csv} -- Pass 2 results
  \item \texttt{data/stage2/owner\_locations\_all.csv} -- merged master file (input to Phase 4)
\end{itemize}

% ============================================================
\section{Phase 3 -- AI Adoption Detection}
% ============================================================

\subsection{Design: Layered Three-Pass Pipeline}

Passes are applied in order of increasing API cost. A repo classified by an earlier
pass is never re-examined. This ensures the most expensive checks run only on the
smallest subset of repos.

\subsection{Pass 1 -- Configuration File Check}

Every repo checked for 12 known agent config files via batch GraphQL (10 repos/request).
These are hard positives -- nobody accidentally creates a \texttt{CLAUDE.md}.

\begin{table}[h]
\centering
\begin{tabular}{@{}ll@{}}
\toprule
File & Agent \\
\midrule
\texttt{CLAUDE.md}                       & Claude Code \\
\texttt{AGENTS.md}                       & Claude Code / Generic \\
\texttt{.cursorrules}, \texttt{.cursorignore} & Cursor \\
\texttt{.aider.conf.yml}                 & Aider \\
\texttt{.roomodes}                       & Roo Code \\
\texttt{.windsurfrules}                  & Windsurf \\
\texttt{GEMINI.md}                       & Gemini CLI \\
\texttt{.github/copilot-instructions.md} & GitHub Copilot \\
\texttt{.continue/config.json}           & Continue \\
\texttt{.codeium/config.json}            & Codeium \\
\texttt{AGENTS.json}                     & OpenAI Codex \\
\bottomrule
\end{tabular}
\end{table}

Any file found: repo immediately labelled \textbf{agentic}. Skips Pass 2 and 3.

\subsection{Pass 2 -- Commit Trailer and PR Branch Scan}

100 most recent commits and 50 most recent PRs fetched per repo. Commit messages
scanned for \texttt{Co-authored-by:} trailer patterns and author email patterns.
PR head branch names scanned for agent prefixes (\texttt{copilot/}, \texttt{devin/},
\texttt{codex/}). Agents covered: Claude Code, GitHub Copilot, Devin, Cursor, Aider,
Windsurf, Gemini, Codeium.

\subsection{Pass 3 -- Deep Scan}

Runs only on repos unclassified after Pass 1 and 2. Up to 300 commits, 100 PRs (body,
labels, branch), and \texttt{package.json} dependency file. npm SDK deps checked:
\texttt{@anthropic-ai/sdk}, \texttt{@anthropic-ai/claude-code}, \texttt{@cursor/agent}.

\subsection{Detection Approach}

Detection is \textbf{heuristic-based only} -- pattern matching against known signatures.
No ML classifier is used because no labeled training dataset exists for the JS/TS/npm
ecosystem. Commit size and behavioral features are not used.

\subsection{Classification Rule (all passes)}

\begin{table}[h]
\centering
\begin{tabular}{@{}ll@{}}
\toprule
Condition & Label \\
\midrule
$\geq$3 signals OR $\geq$2 distinct agents & \textbf{agentic} \\
1--2 signals from a single agent           & \textbf{ai\_assisted} \\
No signal found                            & \textbf{none} \\
\bottomrule
\end{tabular}
\end{table}

\subsection{Single vs Multi-Signal (RQ2)}

Pass 1 config-file-only labels are saved separately alongside the full three-pass labels.
McNemar's test is applied to the paired per-repo outcomes to test whether multi-signal
detection captures significantly more adoption than single-signal.

\subsection{Output Files}
\begin{itemize}
  \item \texttt{data/stage3/pass1\_results.csv} -- config file results
  \item \texttt{data/stage3/pass2\_results.csv} -- commit trailer results
  \item \texttt{data/stage3/pass3\_results.csv} -- deep scan results
  \item \texttt{data/stage3/corpus\_detected\_full.csv} -- final merged labels for all 60,129 repos
\end{itemize}

% ============================================================
\section{Phase 4 -- Geographic Analysis}
% ============================================================

\subsection{Join}

\texttt{corpus\_detected\_full.csv} is joined with \texttt{owner\_locations\_all.csv}
on \texttt{owner\_login}. Each repo inherits its owner's country. Repos with
\texttt{not\_specified} owners are kept in global totals but excluded from country
ranking.

\subsection{Country Aggregation}

For each country with $\geq$30 repos in the corpus:

\begin{itemize}
  \item Total repos
  \item Adopted repos (agentic + ai\_assisted)
  \item Agentic-only repos
  \item Adoption rate (\%) + Wilson 95\% CI
  \item Agentic rate (\%)
\end{itemize}

Countries ranked by adoption rate (descending) to produce the primary result table.

\subsection{Statistical Tests}

\begin{itemize}
  \item \textbf{RQ1} -- One-sample binomial test: global adoption rate $>$ 5\% baseline.
        Null: adoption is trivial (casual one-time use only).

  \item \textbf{RQ2} -- McNemar test on paired repo labels (Pass 1 only vs.\ full
        pipeline). Measures undercounting delta in percentage points.

  \item \textbf{RQ3} -- Chi-squared test of independence across all eligible countries.
        Followed by pairwise Fisher's exact tests (each country vs.\ global baseline)
        with Benjamini-Hochberg FDR correction at $\alpha = 0.05$.
\end{itemize}

\subsection{Output Files}
\begin{itemize}
  \item \texttt{data/stage4/corpus\_with\_country.csv} -- full corpus with country column
  \item \texttt{data/stage4/country\_adoption.csv} -- per-country ranking table
  \item \texttt{data/stage4/statistical\_tests.txt} -- RQ1/RQ2/RQ3 test results
  \item \texttt{data/stage4/pairwise\_fisher.csv} -- pairwise country vs.\ global results
\end{itemize}

% ============================================================
\section{Supplementary: Egypt Corpus (Previous Work)}
% ============================================================

A targeted Egyptian developer corpus was built in an earlier version of this study.
It is retained as supplementary material and is not part of the main geographic ranking.
Key findings from this work are documented for potential use in the discussion chapter.

\begin{table}[h]
\centering
\begin{tabular}{@{}lll@{}}
\toprule
Metric & Version 1 (wrong filter) & Version 2 (corrected) \\
\midrule
Filter used            & pushed $\geq$ 2025-01-01    & created $\geq$ 2025-02-24 \\
Corpus size            & 136 repos                   & 83 repos \\
Multi-signal adoption  & 17.6\% (24/136)             & 41.0\% (34/83) \\
Single-signal adoption & 11.8\% (16/136)             & 25.3\% (21/83) \\
Undercounting delta    & 5.9 pp                      & 15.7 pp \\
\bottomrule
\end{tabular}
\end{table}

The jump from 17.6\% to 41.0\% is not real growth: it reflects the removal of
pre-AI-era repos that were active recently but created before 2025-02-24. This version
comparison is itself a methodological finding illustrating how a single corpus
construction decision can distort adoption measurements by more than 20 percentage points.

Egypt will appear as one data point in the Phase 4 country ranking if enough repos
qualify under the $\geq$30 threshold.

\end{document}
